% !TeX root = ../main.tex
\section{Polynomial IOPs and the IOP-to-SNARK Paradigm}
\label{sec:iop-paradigm}

\subsection{The PCP theorem and local verification}

For functions $r,q:\mathbb{N}\to\mathbb{N}$, the class $\mathrm{PCP}(r(n),q(n))$ contains languages decided by polynomial-time probabilistic verifiers that use at most $r(n)$ random bits and query at most $q(n)$ symbols of a proof oracle. The identity
\[
\mathrm{NP}=\mathrm{PCP}(0,\operatorname{poly}(n))
\]
is essentially the definition of NP: one direction reads an ordinary polynomial-size witness in full, while the other records the polynomially many fixed positions read by a deterministic oracle verifier and uses their contents as an NP witness.

The PCP theorem is the much stronger statement
\begin{tcolorbox}
\[
\mathrm{NP}=\mathrm{PCP}(O(\log n),O(1)).
\]
\end{tcolorbox}
Completeness can be made $1$, while a constant soundness gap can be amplified by repetition. The nontrivial direction constructs a redundant encoding of an NP witness in which a constant number of local consistency checks detect any globally inconsistent proof with constant probability \cite{AroraSafra1998,AroraLundMotwaniSudanSzegedy1998}.

This local-to-global principle is the information-theoretic ancestor of modern IOPs. Kilian's argument makes the cryptographic bridge explicit: commit to a PCP oracle with a Merkle tree, reveal only the verifier's queried symbols and authentication paths, and thereby obtain succinct communication from a locally checkable proof \cite{Kilian1992}. IOPs retain the oracle viewpoint but allow several rounds of oracle messages and challenges; polynomial IOPs further exploit algebraic structure in those oracles.

\subsection{Interactive proofs, IOPs, and polynomial IOPs}

To use polynomial commitments inside proof systems, recall the following hierarchy:
\begin{itemize}[leftmargin=2em]
    \item an \textbf{interactive proof (IP)} exchanges a small number of messages;
    \item an \textbf{interactive oracle proof (IOP)} allows the prover to send long oracles that the verifier only queries in a few places;
    \item a \textbf{polynomial IOP} is an IOP in which the prover's main oracle messages are polynomials.
\end{itemize}

A polynomial IOP for circuit satisfiability has the form
\[
\Setup(C)\to (\pp,\vp),
\]
then the prover sends oracle access to a few polynomials $f_1,\dots,f_t$, the verifier samples random challenges $r_1,\dots,r_s$, and finally checks algebraic identities involving evaluations of those polynomials.

The key point is that the verifier never needs to read entire witness polynomials. It suffices to query them at a few random points.

\subsection{From a polynomial IOP to a SNARK}

Suppose a polynomial IOP commits to $t$ polynomials and the verifier needs $q$ evaluation proofs. If we instantiate the commitment layer with a PCS, then:
\begin{itemize}[leftmargin=2em]
    \item the prover sends $t$ polynomial commitments;
    \item each queried evaluation is accompanied by a PCS opening proof;
    \item the IOP's verifier logic is run on the opened values.
\end{itemize}

If the IOP is public coin, the Fiat--Shamir transform removes interaction by deriving verifier challenges from a hash of the transcript \cite{FiatShamir1986}:
\[
 r_1 = H(C,x,\mathsf{com}_1,\dots),\qquad
 r_2 = H(C,x,\mathsf{com}_1,\dots,r_1,\dots),\quad \text{etc.}
\]

This yields a non-interactive argument in the random-oracle model. Non-interactive zero knowledge in the common-reference-string model originates with Blum, Feldman, and Micali \cite{BlumFeldmanMicali1988}; Fiat--Shamir is a different route, usually modeled through a random oracle.

\subsection{Complexity bookkeeping}

If the PCS commitment size is $\mathsf{pcsCom}$ and the opening size is $\mathsf{pcsOpen}$, then the resulting proof length is roughly
\[
 t\cdot \mathsf{pcsCom} + q\cdot \mathsf{pcsOpen}.
\]
The verifier's work is roughly
\[
 q\cdot \mathrm{time}(\mathsf{PCS\ verify}) + \mathrm{time}(\text{IOP verifier}),
\]
and the prover's work is roughly
\[
 t\cdot \mathrm{time}(\mathsf{Commit}) + q\cdot \mathrm{time}(\mathsf{Open}) + \mathrm{time}(\text{IOP prover}).
\]

This formula explains why the same polynomial IOP can lead to systems with different tradeoffs:
\begin{center}
\small
\begin{tabular}{@{}p{0.19\textwidth}p{0.29\textwidth}p{0.42\textwidth}@{}}
\toprule
PCS & Resulting flavor & Typical tradeoff \\
\midrule
KZG \cite{KateZaveruchaGoldberg2010} & pairing-based Plonk & short proofs; structured setup \\
IPA \cite{Bulletproofs2018} & IPA-based Plonkish systems & transparent setup; slower verification \\
FRI \cite{BenSassonBentovHoreshRiabzev2018FRI} & STARK-like systems & transparent and hash-based; larger proofs \\
\bottomrule
\end{tabular}
\end{center}
